% 【本文件写什么】api-v0 契约层，叶子，只写语义不写实现
% - crate 路径：os/components/wateros-driver/driver-network/network-api/api-v0/src/
% - 列出并说明：pub trait、核心类型、错误枚举、常量
% - 每个公开 API 的前置条件、后置条件、线程/中断上下文限制
% - 与 Linux/用户态约定的对齐点与 deliberate 差异
% - v0 版本当前行为与后续可替换 extension 点
% 【事实来源】
% - 上述 crate 下全部 .rs 源文件
% - docs/exports/public-api/ 中对应符号表，以源码为准
% 【不写什么】
% - impl 内部数据结构、汇编、具体算法步骤

\subsubsection{network-api-v0}

\paragraph{NetworkDevice契约。}设备须实现\code{mac\_address}、\code{send}与\code{receive}；\code{send}/\code{receive}操作完整以太网L2帧，含目的/源MAC与EtherType。\code{mtu}默认\code{DEFAULT\_MTU==1500}；\code{is\_link\_up}默认\code{true}，占位设备可覆盖为\code{false}。\code{receive}在\code{buf}不足以容纳帧时返回\code{InvalidParam}；无数据时允许返回0字节。

\paragraph{全局注册表。}与块子系统相同，\code{register\_network\_device}返回稳定下标；\code{first\_network\_device}供smoltcp\code{init}选取首卡。多网卡场景表可存多实例，但协议栈v0仅绑定首张注册网卡。

\begin{rustcode}
pub const DEFAULT_MTU: usize = 1500;

pub trait NetworkDevice: Send {
    fn mac_address(&self) -> [u8; 6];
    fn mtu(&self) -> usize { DEFAULT_MTU }
    fn is_link_up(&self) -> bool { true }
    fn send(&mut self, buf: &[u8]) -> DriverResult<()>;
    fn receive(&mut self, buf: &mut [u8]) -> DriverResult<usize>;
}
\end{rustcode}

\paragraph{与协议栈边界。}本crate不定义IP/socket语义；\code{impl-smoltcp}通过\code{SmoltcpAdapter}实现\code{smoltcp::phy::Device}，在token生命周期内完成帧搬运。错误仍用\code{DriverError}，socket层使用\code{\&'static str}短消息，见 \code{stack} 模块。
